Recent work suggests that Large Language Models (LLMs) can learn novel geometric structures in-context from symbolic data. However, it remains unclear whether the ability to retrieve specific components of these geometries—which we term \addressing—implies an underlying spatial understanding. In this work, we investigate the relationship between \addressing and \reasoning in \gpt using the \kilogram dataset of complex 2D tangram shapes. We find that \gpt achieves perfect \addressing accuracy (100\%) even with zero or minimal repetition, demonstrating a high capacity for symbolic mapping. In contrast, its performance on \reasoning tasks (e.g., relative position queries) remains significantly lower, ranging from 50\% to 80\%, and shows only gradual improvement with up to 20 examples. Our results reveal a fundamental decoupling between symbolic lookup and spatial perception: while models can accurately "address" the coordinates of a named part, they struggle to "perceive" the geometric relationships without specialized processing. This gap highlights a limitation in current in-context learning paradigms for spatial tasks.
